% listings package configuration
\usepackage{listings}
\usepackage{fvextra}
\usepackage{luacolor}
\usepackage{lua-ul}
%                         https://tex.stackexchange.com/a/466224/151207

%\usepackage{xcolor} % Required for color definitions
% see includes-colours.tex

% Define custom colors
\definecolor{keywordcolor}{rgb}{0.13,0.29,0.53}
\definecolor{stringcolor}{rgb}{0.31,0.60,0.02}
\definecolor{commentcolor}{rgb}{0.56,0.35,0.01}
\definecolor{backcolor}{rgb}{0.95,0.95,0.92}

\definecolor{codecolor}{rgb}{0.94, 0.92, 0.84} % default code background colour
\definecolor{codeframe}{rgb}{0.84, 0.79, 0.69} % default code frame colour

% Julia language definition
\lstdefinelanguage{Julia}{
	basicstyle=\linespread{3}\fontsize{7}{9}\selectfont\fontspec{Fira Code},
	keywordsprefix=\@,
	morekeywords={
			exit,whos,edit,load,is,isa,isequal,typeof,tuple,ntuple,uid,hash,finalizer,convert,promote,
			subtype,typemin,typemax,realmin,realmax,sizeof,eps,promote_type,method_exists,applicable,
			invoke,dlopen,dlsym,system,error,throw,assert,new,Inf,Nan,pi,im,begin,while,for,in,return,
			break,continue,macro,quote,let,if,elseif,else,try,catch,end,bitstype,ccall,do,using,module,
			import,export,importall,baremodule,immutable,local,global,const,Bool,Int,Int8,Int16,Int32,
			Int64,Uint,Uint8,Uint16,Uint32,Uint64,Float32,Float64,Complex64,Complex128,Any,Nothing,None,
			function,type,typealias,abstract},
	sensitive=true,
	alsoother={$},
	morecomment=[l]{\#},
	morecomment=[s]{/*}{*/},
	morestring=[b]",
	morestring=[b]',
}

% General listings settings
\lstset{% 
	basicstyle=\linespread{1.3}\footnotesize\ttfamily,      % font size, mono font, 
	% linespread fixes white tripes
	backgroundcolor=\color{codecolor},  % background color
	showspaces=false,              % show spaces (with underscores)
	frame=single,                  % adds a frame around the code
	tabsize=2,                     % default tabsize
	breaklines=true,               % automatic line breaking
	columns=fullflexible,
	framerule=0.5pt,
	rulecolor=\color{codeframe},
	framesep=1em,                  % padding
	xleftmargin=1em,
	columns=fullflexible,          % make sure to use fixed-width font, CM typewriter is NOT fixed width
	numbers=none,                  % left
	numberstyle=\footnotesize\ttfamily\color{codeframe},
	numbersep=5pt,                 % how far the line numbers are from the code
	stepnumber=1,
	numberfirstline=true,
	numberblanklines=true,
	tabsize=4,
	lineskip=-1.5pt,
	extendedchars=true,
	mathescape=false,
	breaklines=true,
	keywordstyle=\color{keywordcolor}\bfseries,
	identifierstyle=, % using emph or index keywords
	commentstyle=\color{commentcolor},
	stringstyle=\color{stringcolor},
	showstringspaces=false,
	showtabs=false,
	upquote=false,
	%texcl=true, % interpet comments as LaTeX
	literate={_}{\_}1
}

% Robust inline code.  This uses fvextra's verbatim reader rather than wrapping
% \lstinline in another macro, so raw underscores and long paths are safe.
% commandchars keeps legacy notes such as \code{foo\_bar} readable.
% Paper style: inline code as plain typewriter, no highlight background
\renewcommand{\FancyVerbFormatInline}[1]{#1}
\newcommand{\code}{\Verb[
	breaklines=true,
	breakanywhere=true,
	breaksymbolleft={},
	breaksymbolright={},
	formatcom=\codeblockfont\small,
	commandchars=\\\{\}
]}
\pdfstringdefDisableCommands{%
	\def\code#1{#1}%
}

% Plain verbatim blocks. Keep this on FancyVerb/fvextra rather than listings:
% listings inline state can leak into aux/toc writes in recent TeX Live.
\makeatletter
\define@key{FV}{plancolor}[]{\fvset{bgcolor=plancolor,fillcolor=\color{plancolor}}}
\define@key{FV}{robocupcolor}[]{\fvset{bgcolor=robocupcolor,fillcolor=\color{robocupcolor}}}
\define@key{FV}{sepiacolor}[]{\fvset{bgcolor=sepiacolor,fillcolor=\color{sepiacolor}}}
\define@key{FV}{donecolor}[]{\fvset{bgcolor=donecolor,fillcolor=\color{donecolor}}}
\makeatother
\DefineVerbatimEnvironment{lstplain}{Verbatim}{
	fontsize=\scriptsize,
	baselinestretch=0.95,
	formatcom=\codeblockfont,
	breaklines=true,
	breakanywhere=true,
	breaksymbolleft={},
	breaksymbolright={},
	bgcolor=codecolor,
	frame=single,
	framerule=0.5pt,
	framesep=1em,
	xleftmargin=0em,
	rulecolor=\color{codeframe},
	fillcolor=\color{codecolor}
}

\DefineVerbatimEnvironment{vrb}{Verbatim}{
	fontsize=\scriptsize,
	baselinestretch=0.95,
	formatcom=\codeblockfont,
	breaklines=true,
	breakanywhere=true,
	breaksymbolleft={},
	breaksymbolright={},
	bgcolor=codecolor,
	frame=single,
	framerule=0.5pt,
	framesep=1em,
	xleftmargin=0em,
	rulecolor=\color{codeframe},
	fillcolor=\color{codecolor}
}


\lstnewenvironment{lstjulia}[1][codecolor]{
	\lstset{% 
		language=Julia,
		basicstyle=\fontsize{7}{9}\selectfont\fontspec{Fira Code},
		xleftmargin=0em,
		backgroundcolor=\color{#1},
		commentstyle=\color{commentcolor},
		keywordstyle=\color{keywordcolor},
		stringstyle=\color{stringcolor},
		numberstyle=\tiny\color{gray},
		breaklines=true,
		showstringspaces=false,
		frame=single,
		tabsize=2,
		literate={_}{\_}1
	}
}{}

\lstnewenvironment{lstpython}[1][codecolor]{
	\lstset{%
		language=Python,
		basicstyle=\fontsize{7}{9}\selectfont\fontspec{Fira Code},
		xleftmargin=0em,
		backgroundcolor=\color{#1},
		commentstyle=\color{commentcolor},
		keywordstyle=\color{keywordcolor},
		stringstyle=\color{stringcolor},
		numberstyle=\tiny\color{gray},
		breaklines=true,
		showstringspaces=false,
		frame=single,
		tabsize=2,
		literate={_}{\_}1
	}
}{}
